\clearpage
\hypertarget{native-473-lower}{}
\begin{center}
\rule{0.55\linewidth}{0.4pt}\\[2pt]
Astronomische Nachrichten. Band XXII. Nr. 516. S. 189..192. 3. Oktober 1844.\\[2pt]
\rule{0.55\linewidth}{0.4pt}
\end{center}
\begin{flushright}1. September 1844.\end{flushright}

Gestern Abend konnte ich nach langer Unterbrechung wieder ein paar Vergleiche des Kometen mit benachbarten Fixsternen machen. Sie eigneten sich aber nur zur Bestimmung der Rektaszension

\begin{center}
1844 August 31.\quad $8^{\rm u}\ 25^{\rm m}\ 57^{\rm s}$ m. Z.\quad Rektasz. $203^\circ\ 45'\ 8^{\rm u}3$
\end{center}

Die Parabel von Dr. \textsc{Goldschmidt} stimmt damit noch bis auf eine Sekunde überein.

Ich merke noch an, dass bei \textsc{Goldschmidts} Formeln zur Berechnung der Koordinaten, wie sie in Nr. 512 der A. N. abgedruckt sind, noch der Radiusvektor als Faktor hinzukommen muss.

\begin{flushright}Göttingen, Sept. 3.\end{flushright}

Ich schicke Ihnen noch eine Kometenbeobachtung.

\begin{center}
1844 Sept. 2.\quad $8^{\rm u}\ 27^{\rm m}\ 11^{\rm s}$\quad Rektasz. $=203^\circ\ 11'\ 0''2$\\[2pt]
Dekl. $=+7^\circ\ 6'\ 32''1$
\end{center}

Sie beruht auf zwei Vergleichen mit einem Stern 9. Größe aus \textsc{Bessels} Zonen 83 und 160 und mit den beiden \textsc{Piazzi}schen Sternen XIII, 194 und 208. --- Die Rektaszension des zweiten Sterns (XIII, 194) hat bei mir einige Bedenken geweckt, und es scheint fast, als hätte der Stern eine Eigenbewegung. Die Zeitdifferenz zwischen dem zweiten und dritten Stern ist wie folgt angegeben

\begin{center}
\setlength{\tabcolsep}{3pt}
\begin{tabular}{lll}
in der Hist. Céleste & $2^{\rm m}\ 26^{\rm s}00$ & \\
in \textsc{Zachs} Katalog & $2\ 26.11$ & \\
in \textsc{Piazzis} Katalog, 1. Aufl. & $2\ 26.04$ & jeder Stern 4-mal \\
\hphantom{in \textsc{Piazzis} Katalog, }2. Aufl. & $2\ 26.31$ & jeder Stern 10-mal beob. \\
in \textsc{Bessels} Zone 160. & $2\ 27.00$ & \\
meine beiden Durchgänge ergeben & $2\ 27.64$ & \\
\end{tabular}
\end{center}

\Needspace{6\baselineskip}
\SourceLeaf{474}
Die Differenz scheint also zuzunehmen, obwohl die jährliche Präzession für XIII. 208 um ein paar Hundertstel einer Bogensekunde kleiner ist als für XIII, 194. --- Für die obige Positionsbestimmung des Kometen habe ich nur die \textsc{Bessel}schen Beobachtungen der drei Sterne verwendet und das Mittel aus allen Ergebnissen angesetzt. Würde ich bei den Rektaszensionen das Ergebnis aus XIII, 194 ausschließen, ergäbe sich
\[
203^\circ\ 10'\ 57''2
\]
\noindent Bei der Reduktion ist die Refraktion übrigens genau berücksichtigt. --- \textsc{Goldschmidts} Parabel gibt die Rektaszension noch gut wieder, ergibt aber eine um 1 Minute zu kleine Deklination.

\begin{flushright}Göttingen, Sept. 6.\end{flushright}

Auch gestern habe ich den Kometen beobachtet. Ich konnte aber nur einen Durchgang erfassen, der auch nur für die \emph{Deklination} brauchbar war. Das unter Berücksichtigung der Rektaszension beobachtete Ergebnis war

\begin{center}
1844 Sept. 5.\quad $7^{\rm u}\ 56^{\rm m}\ 56^{\rm s}$ m. Z.\quad Dekl. des Kometen $+5^\circ\ 21'\ 20''2$
\end{center}

\textsc{Goldschmidts} Parabel ergibt $1'\ 5''$ weniger.

\begin{flushright}Göttingen, Sept. 8.\end{flushright}

Heute schicke ich Ihnen meine Kometenbeobachtung von vorgestern:

\begin{center}
\setlength{\tabcolsep}{3pt}
\begin{tabular}{lll}
1844 Sept. 6. & $8^{\rm u}\ 1^{\rm m}\ 53^{\rm s}$ & Rektaszension $202^\circ\ 7'\ 6''7$ \\
 & $8\ 8\ 7$ & Nördl. Deklination $4\ 46\ 25.6$ \\
\end{tabular}
\end{center}

Die Rektaszension beruht auf vier, die Deklination auf zwei Vergleichen.

Gestern Abend war der Himmel bedeckt.
\begin{flushright}C. F. \textsc{Gauss}.\end{flushright}
